\section{Group Presentations}

\exercise{40.1}{Presentations of $\Z_4$}
One, two, and three generators can be used as follows:
\[
\gen{a\mid a^4=1},\qquad
\gen{a,b\mid a^4=1,\ b=a^2},
\]
\[
\gen{a,b,c\mid a^4=1,\ b=a^2,\ c=a^3}.
\]
In each case eliminate the additional generators using their defining relations. Every element is a power of $a$, with exponent reduced modulo 4, so there are at most four elements. Sending $a$ to $1\in\Z_4$, $b$ to 2 and $c$ to 3 (where present) satisfies the relations and maps onto $\Z_4$, proving there are exactly four elements.

\exercise{40.2}{A three-generator presentation of $S_3$}
Use
\[
\gen{a,b,c\mid a^3=1,\ b^2=1,\ bab=a^{-1},\ c=ab}.
\]
The relation for $c$ lets us eliminate it. The conjugation relation gives $ba=a^{-1}b$, so every word becomes $a^ib^j$ with $0\le i<3$, $0\le j<2$. Thus there are at most six elements.
Map $a\mapsto(123)$, $b\mapsto(12)$, $c\mapsto(123)(12)$. These permutations satisfy the relations and generate $S_3$, so the presented group has at least six elements. It is therefore isomorphic to $S_3$.

\exercise{40.3}{Tables for the octic and quaternion groups}
\idea Write every element as $a^ib^j$, with $0\le i<4$, $j\in\{0,1\}$. The rule $ba=a^{-1}b$ moves $b$ past a power of $a$. If $b^2=a^c$, then
\[
(a^ib^j)(a^ub^v)=a^{\,i+(-1)^ju+cjv}b^{\,j+v\bmod2},
\]
where the exponent of $a$ is reduced modulo 4. Use $c=0$ for the octic group and $c=2$ for the quaternion group.
The octic table, in the requested order, is
\[\begin{array}{c|cccccccc}\cdot&1&a&a^2&a^3&b&ab&a^2b&a^3b\\\hline
1&1&a&a^2&a^3&b&ab&a^2b&a^3b\\
a&a&a^2&a^3&1&ab&a^2b&a^3b&b\\
a^2&a^2&a^3&1&a&a^2b&a^3b&b&ab\\
a^3&a^3&1&a&a^2&a^3b&b&ab&a^2b\\
b&b&a^3b&a^2b&ab&1&a^3&a^2&a\\
ab&ab&b&a^3b&a^2b&a&1&a^3&a^2\\
a^2b&a^2b&ab&b&a^3b&a^2&a&1&a^3\\
a^3b&a^3b&a^2b&ab&b&a^3&a^2&a&1\\\end{array}\]The quaternion table is
\[\begin{array}{c|cccccccc}\cdot&1&a&a^2&a^3&b&ab&a^2b&a^3b\\\hline
1&1&a&a^2&a^3&b&ab&a^2b&a^3b\\
a&a&a^2&a^3&1&ab&a^2b&a^3b&b\\
a^2&a^2&a^3&1&a&a^2b&a^3b&b&ab\\
a^3&a^3&1&a&a^2&a^3b&b&ab&a^2b\\
b&b&a^3b&a^2b&ab&a^2&a&1&a^3\\
ab&ab&b&a^3b&a^2b&a^3&a^2&a&1\\
a^2b&a^2b&ab&b&a^3b&1&a^3&a^2&a\\
a^3b&a^3b&a^2b&ab&b&a&1&a^3&a^2\\\end{array}\]
For example, the product $(a^2b)(a^3b)$ is $a^{2-3}=a^3$ in the first group and $a^{2-3+2}=a$ in the second. This checks the effect of the different square of $b$.

\exercise{40.4}{All groups of order 14}
Sylow gives a unique order-7 subgroup $\gen a$, so it is normal. Cauchy's theorem gives $b$ of order 2 outside it. The elements $a^ib^j$ exhaust the group, since the subgroup generated by $a,b$ has order divisible by both 7 and 2.
Conjugation has $bab^{-1}=a^r$ with $r$ nonzero modulo 7. Since $b^2=1$, $r^2\equiv1\pmod7$, so $(r-1)(r+1)=0$ in the field $\Z_7$ and $r=1$ or $6$.
The first case gives $\Z_{14}$. The second gives
$\gen{a,b\mid a^7=b^2=1,\ ba=a^6b}$, the dihedral group $D_7$. Exercise 40.13 proves it has order 14; it is nonabelian because inversion acts nontrivially on $a$.
These are the only two possibilities, and their different commutativity makes them nonisomorphic.

\exercise{40.5}{All groups of order 21}
The order-7 Sylow subgroup $\gen a$ is normal. Choose $b$ of order 3 by Cauchy's theorem; it is outside $\gen a$, and together the elements generate all 21 elements.
The action $bab^{-1}=a^r$ must satisfy $r^3\equiv1\pmod7$. Checking the six nonzero residues gives $r=1,2,4$.
For $r=1$ the group is $\Z_{21}$. For $r=2$ we obtain
\[
\gen{a,b\mid a^7=b^3=1,\ ba=a^2b},
\]
a nonabelian group of order 21 by Exercise 40.13. The choice $r=4$ gives the same isomorphism type: replacing $b$ by $b^2$ changes the action exponent from 2 to $2^2=4$, while $b^2$ still generates the order-3 subgroup. Thus there is exactly one abelian and one nonabelian isomorphism type.

\exercise{40.6}{Consequences of relators}
Include conjugates, which are missing from the proposed definition. A consequence is an element of the normal closure of the relators in the free group, equivalently a finite product
\[
\prod_{j=1}^k w_j r_{i_j}^{\epsilon_j}w_j^{-1},
\qquad\epsilon_j\in\{1,-1\}.
\]
Allowing arbitrary integer powers gives the same normal closure. The empty product is the identity. Products merely of powers of the original relators need not form a normal subgroup.

\exercise{40.7}{Isomorphic presentations}
Two presentations are isomorphic when the groups they present are isomorphic. A bijective renaming of generators and relators is sufficient but not necessary.
For example $\gen{a\mid a^4=1}$ and $\gen{a,b\mid a^4=1,b=a^2}$ both present $\Z_4$, as proved in Exercise 40.1, but their generator sets have different sizes, so no bijective renaming exists. The printed proposed definition is therefore too restrictive.

\exercise{40.8}{True or false: presentations}
\begin{enumerate}[label=\textbf{(\alph*)}]
\item True. Take a free group mapping onto the group and use all kernel elements as relators.
\item True. Introduce arbitrarily many new generators together with relations expressing each as an old word.
\item False. By the definition used here, all presentations of the same group are isomorphic presentations.
\item False. A finite presentation gives a finitely generated, hence countable, group. The uncountable additive group $\R$ cannot have one.
\item False. $\gen{a\mid\ }$ presents the infinite cyclic group with a finite generator set and no relators.
\item True. A cyclic group has $\gen{a\mid a^n=1}$ if its order is $n$, or $\gen{a\mid\ }$ if it is infinite. For the trivial group use $n=1$.
\item True. Consequences form the normal closure, which contains every conjugate of each relator.
\item False. $\gen{a\mid a^2=1}$ and $\gen{a\mid a^3=1}$ have the same number of generators but give nonisomorphic groups of orders 2 and 3.
\item True. If $F/R$ is abelian, all commutators in $F$ map to the identity, so $[F,F]\subseteq R$.
\item False. $\gen{a,b\mid b=1}$ presents the free group $\Z$ but has the nonidentity free word $b$ as a relator.
\end{enumerate}

\exercise{40.9}{Groups of order 15 are abelian}
Sylow gives a normal subgroup $\gen a$ of order 5. Cauchy's theorem gives $b$ of order 3, and $a,b$ generate the group. If $bab^{-1}=a^r$, conjugating three times gives $r^3\equiv1\pmod5$.
Among the nonzero residues modulo 5 only $r=1$ satisfies this (or use that their multiplicative group has order 4, coprime to 3). Thus $a$ and $b$ commute.
The group is generated by commuting elements, hence abelian; indeed $ab$ has order 15, so it is cyclic. There is no nonabelian group of order 15.

\exercise{40.10}{An order-six nonabelian presentation}
Apply Exercise 40.13 with $m=3,n=2,r=2$. Since $2^2\equiv1\pmod3$, its normal-form model has exactly $3\cdot2=6$ elements.
It realizes $a$ with order 3 and $ba=a^2b$. If $a,b$ commuted, cancellation would give $a=a^2$, hence $a=1$, contradicting its order 3. Thus the group is nonabelian.

\exercise{40.11}{Uniqueness of the nonabelian group of order six}
For any group of order 6, the Sylow 3-subgroup $\gen a$ is unique and normal. An element $b$ of order 2 exists by Cauchy's theorem and lies outside it, so $a,b$ generate.
Conjugation sends $a$ to $a$ or $a^2$, the two generators of $\gen a$. The first case makes the group abelian and cyclic of order 6. The second gives exactly the presentation in Exercise 40.10, which has order 6 and thus maps isomorphically onto the given group.
Since $S_3$ is nonabelian of order 6, it realizes the second case. Concretely one may send $a$ to $(123)$ and $b$ to $(12)$.

\exercise{40.12}{No subgroup of order six in $A_4$}
An order-six subgroup would, by Exercise 40.11, be cyclic of order 6 or isomorphic to $S_3$.
The cyclic possibility is impossible because elements of $A_4$ have orders only $1,2,3$: they are the identity, double transpositions, and 3-cycles.
The $S_3$ possibility is also impossible. Its three involutions do not commute, and the product of two distinct involutions has order 3. In $A_4$, the only involutions are the three double transpositions; they form a Klein four group, commute, and multiply in pairs to another involution. An isomorphism preserves these products and orders, giving a contradiction.

\exercise{40.13}{The normal-form multiplication criterion}
\textbf{(a)} Write a formal element as $(s,t)$ with $s\in\Z_m$ and $0\le t<n$, and define
\[
(s,t)(u,v)=(s+r^tu\bmod m,\ t+v\bmod n).
\]
Assume $r^n\equiv1\pmod m$. Then $r$ is invertible modulo $m$, and replacing $t+v$ by its residue modulo $n$ leaves $r^{t+v}$ unchanged modulo $m$. Thus the first coordinates of both associations of a triple product equal
$s+r^tu+r^{t+v}w$, and the second equal $t+v+z\pmod n$. The operation is associative.
The identity is $(0,0)$. For $(s,t)$ take $v=-t\pmod n$ and $u=-r^{-t}s\pmod m$; the resulting $(u,v)$ is a two-sided inverse. This proves a group of order $mn$.
Conversely, for $n>1$ suppose the operation is a group. Then $a=(1,0)$ has order $m$, $b=(0,1)$ has order $n$, and $bab^{-1}=a^r$. Repeating conjugation $n$ times gives $a=a^{r^n}$, hence $r^n\equiv1\pmod m$.
\emph{Degenerate case:} if $n=1$, the displayed operation is always just addition in $\Z_m$, independent of $r$. Thus the literal ``if and only if'' in part (a) needs $n>1$; its sufficiency remains valid for $n=1$.
\textbf{(b)} If the congruence holds, the group constructed above satisfies the presentation relations. Conversely those relations put every word in one of the $mn$ forms $a^ib^j$, so the presented group maps onto this order-$mn$ model and has at most $mn$ elements; it therefore has exactly $mn$.
If the presentation has order $mn$, its $mn$ normal forms must all be distinct. In particular $a$ has order $m$. The relations $b^n=1$ and $bab^{-1}=a^r$ imply $a=a^{r^n}$ and hence the congruence. This also proves part (b) for $n=1$, where $b=1$ forces $a=a^r$.

\exercise{40.14}{A unique nonabelian group of order $pq$}
Let $p<q$ be primes with $p\mid q-1$. For any order-$pq$ group, the Sylow $q$-subgroup is unique because its number divides $p<q$ and is 1 modulo $q$. Write it $\gen a$, and choose $b$ of order $p$ by Cauchy's theorem. The two subgroups intersect trivially and generate the group, so it has a presentation
\[
\gen{a,b\mid a^q=b^p=1,\ ba=a^rb},\qquad r^p\equiv1\pmod q.
\]
The unit group $\Z_q^*$ is cyclic of order $q-1$, so the $p$th roots of 1 form its unique order-$p$ subgroup. Its nonidentity elements are $r,r^2,\ldots,r^{p-1}$ for any one generator $r$.
The root 1 gives the abelian cyclic group. Every other root gives a nonabelian order-$pq$ group by Exercise 40.13. Replacing $b$ by $b^j$, $1\le j<p$, changes the action exponent to $r^j$ without changing its generated order-$p$ subgroup. Therefore all nontrivial roots give isomorphic presentations, proving existence and uniqueness of the nonabelian isomorphism type.
